% reality/topology_and_special_elements.tex

\chapter{Topology and Special Elements}

\label{chap:topology-and-special-elements}

Railway infrastructure is networked. Reality interpretation therefore
requires explicit topology context in addition to geometric state.

% ============================================================
\section{Topology Context}
% ============================================================

\begin{definition}[Railway network context]
Railway network context is the set of connectivity relations between
aligned track segments, branches, merges, and route continuations.
\end{definition}

Topology context governs route continuity and interpretation of station,
projection, and observation data in multi-track systems.

% ============================================================
\section{Special Elements}
% ============================================================

Switches, crossings, station throats, and yard structures are special
elements where geometric and operational branching is explicit.

\begin{definition}[Connection point]
A connection point is a location where two or more alignment segments
are topologically linked.
\end{definition}

These elements do not redefine constructive alignment identity; they
provide network context in which identities are interpreted.

% ============================================================
\section{Geometry--Topology Separation}
% ============================================================

A useful abstraction separates geometry from connectivity while allowing
explicit coupling at interpretation time.

\begin{definition}[Graph representation]
A railway network may be represented as a graph in which nodes denote
connection points and edges denote alignment segments.
\end{definition}

This representation clarifies that topology context and geometric state
are distinct but jointly required in Reality workflows.

% ============================================================
\section{Data Ambiguity and Inference}
% ============================================================

Observed datasets frequently under-specify topology. Inference of route
continuity and branch assignment is therefore part of reality
interpretation and must be handled explicitly.

% ============================================================
\section{Connection to AIM}
% ============================================================

\begin{summary}
Topology and special elements provide the network context required for
real-world interpretation of measured and realized railway states.

Within AIM, topology is consumed as contextual structure and does not
replace constructive alignment semantics.
\end{summary}
